\section{The regime that answers: freezing the field}
\label{sec:ssrhalf}

Every value of the dial so far is an assumption.  This section builds
the one transport in the book that is \emph{derived}.  The ingredient
is Chapter~4's local-volatility field: one function $v(y)$ of the
normalized strike $y=K/F$ --- a local variance at every level the
underlying might visit --- chosen so that the model reprices every
expiry of the surface at once.  (Chapter~4 wrote $v(\tau,y)$; the
fields of this section and the next are deliberately
time-homogeneous, so $v$ depends on the level alone.)  Now make one
hypothesis: \emph{the field is a property of the underlying at
absolute price levels}.  If the local volatility at $\$650$ is high,
it is high because of what the market believes about the world at
$\$650$ --- not because of where spot happens to stand today.  Under
that hypothesis, when spot moves, the field \emph{stays where it is}
in absolute strike, and the new smile is whatever the frozen field
reprices to.  (Mechanically: after the forward moves to
$F_{\mathrm{old}}e^{H}$, the same absolute strikes carry the same
local variances under new labels, so the normalized coefficient is
read as $v(y\,e^{H})$ --- the quote rule of \cref{def:ssrsigns},
applied to the field.)  No dial is left: the transport is a
computation.

The computation has a famous shortcut, and the shortcut explains the
number $2$ that has been hanging over the chapter since
\cref{sec:ssrdial}.  For short maturities, an option's implied
volatility is approximately an \emph{average of the local volatility
along the way from the money to the strike}: the price is built from
paths that drift from spot toward the strike, sampling the field as
they go.  (Made precise, the right average is a \emph{harmonic} one --- the
reciprocal of the averaged reciprocal --- along the connecting chord,
a result of \citet{BerestyckiBuscaFlorent2002}; the midpoint reading
used below goes back to \citet{HaganKumarLesniewskiWoodward2002}.
The distinctions do not matter here: to first order in the field's
slope every such average agrees, and first order is all the half-rule
needs.)  Write the field in
volatility form and in log-moneyness, $\sqrt{v}$ as a function of
$x=\log y$, and take it to be a straight line with slope $b$ near the
money.  Then the chord average is elementary:
\[
\sigma(k)\;\approx\;\frac{1}{k}\int_{0}^{k}
\big(\sigma_{0}+b\,x\big)\,\dd x
=\sigma_{0}+\frac{b}{2}\,k
\qquad\Longrightarrow\qquad
s_0=\frac{b}{2}.
\]
This is the \textbf{half-rule}: \emph{the implied skew is half the
field's slope}.  The market's smile shows half of what the generator
carries, because each strike only averages the field part of the way
out.

\begin{figure}[htbp]
\centering
\paperfig{\textwidth}{fig_ssr_half.pdf}
\caption{The frozen field's answer (synthetic straight-line field;
panel~(c) the frozen hero smile).  (a)~The generator (dashed, slope
\MacSsrHalfSlopeLoc) and the implied smile it produces at
$\tau=0.10$ (blue, measured slope \MacSsrHalfSlopeImp): the half-rule,
measured ratio \MacSsrHalfRatio.  (b)~The field frozen in absolute
strike, the forward moved $-4\%$: the implied ATM rises
\MacSsrHalfAtmMoveBp\ vol bp against the half-rule's prediction
$2s_0H=\MacSsrHalfAtmPredBp$ vol bp.  (c)~The two closed-form
transports on the real node at $H=-\MacSsrHalfMoveWingPct\%$: midpoint
relabeling minus linear $\mathcal{R}=2$ transport, zero at the money,
$+\MacSsrHalfGapPutBp$/$-\MacSsrHalfGapCallBp$ vol bp in the wings.}
\label{fig:ssrhalf}
\end{figure}

The rule is not folklore; it is measurable.  \Cref{fig:ssrhalf}(a)
shows a synthetic world built for exactly this test: a
time-homogeneous field whose volatility is a straight line in
log-strike with slope \MacSsrHalfSlopeLoc\ (dashed), and the implied
smile obtained by actually pricing through the forward equation with
Chapter~4's marching scheme and inverting Black (blue), at maturity
$\tau=0.10$.  The measured implied slope is \MacSsrHalfSlopeImp\ ---
ratio \MacSsrHalfRatio\ to the generator.  Half, to a tenth of a
percent, with no fitting anywhere.

Now freeze and move.  Let the forward move by $H$ --- take $H<0$, a
drop --- and hold the field.  At the \emph{new} money, the implied
volatility averages the field around the new spot, so it reads the
straight generator there:
\[
\sigma_{\mathrm{atm}}^{\mathrm{new}}
\approx\sigma_{0}+b\,H
=\sigma_{0}+2\,s_0\,H ,
\]
an \emph{increase} for a drop, since $b$ and $H$ are both negative.
The ATM reading slides along the \emph{generator's} slope $b$ --- and
the smile only ever showed half of $b$.  The frozen field therefore
moves the ATM volatility by \emph{twice the skew} per unit move:
$\mathcal{R}=2$, derived.  The overshoot of
\cref{sec:ssrdial}'s taxonomy is the market paying out the half of the
slope the smile had been holding back.  Panel~(b) runs the experiment
in the synthetic world: forward down $4\%$, field untouched, both
smiles repriced through the forward equation.  The implied curve
slides up its generator; the measured ATM change is
\MacSsrHalfAtmMoveBp\ vol bp against the half-rule's
$2s_0H=\MacSsrHalfAtmPredBp$ vol bp.  Exact to the basis point.

At the money, then, the frozen field's answer is settled.  Away from
the money one can push the same logic further and ask for a closed
form: which \emph{old} smile point should tomorrow's smile at label
$k$ read?  The tool is the \emph{midpoint rule}, and it is one line
from what was just derived: for a straight-line field, the average
along the chord equals the field's value at the chord's
\emph{midpoint}, so the implied volatility at strike $K$ reflects the
field near the dollar midpoint $(F+K)/2$.  Two (forward, strike)
pairs that share a midpoint therefore share a reading.
Tomorrow's pair is $(F_{\mathrm{old}}e^{H},\,K)$ with
$K=F_{\mathrm{old}}e^{H}e^{k}$; the old strike $K''$ sharing its
midpoint solves
\[
\frac{F_{\mathrm{old}}+K''}{2}
=\frac{F_{\mathrm{old}}e^{H}+K}{2}
\qquad\Longrightarrow\qquad
K''=F_{\mathrm{old}}\big(e^{H}(1+e^{k})-1\big),
\]
and its old label is the log of that ratio:
\begin{equation}
\ell(k,H)=\log\!\big(e^{H}(1+e^{k})-1\big),
\qquad
w_{\mathrm{new}}(k)=w_{\mathrm{old}}\big(\ell(k,H)\big)
\label{eq:ssrell}
\end{equation}
--- the \emph{midpoint relabeling}, written in total variance (the
chart Chapter~2 fits in; at a single maturity the volatility form
says the same thing, since $\sigma=\sqrt{w/\tau}$).  Check
it at zero move: $\ell(k,0)=\log(1+e^{k}-1)=k$, the identity.  Expand
it in the move: differentiating,
\[
\frac{\partial\ell}{\partial H}
=\frac{e^{H}(1+e^{k})}{e^{H}(1+e^{k})-1},
\qquad
\frac{\partial\ell}{\partial H}\bigg|_{H=0}
=\frac{1+e^{k}}{e^{k}}
=1+e^{-k},
\]
so
\begin{equation}
\ell(k,H)=k+\big(1+e^{-k}\big)H+O(H^{2}) .
\label{eq:ssrellexp}
\end{equation}
At the money, $\ell(0,H)\approx 2H$: the new smile reads the old curve
\emph{two} units of moneyness per unit of move --- the half-rule's
factor again, now recovered from pure geometry, which is a
reassuring consistency check.  Away from the money the displacement
$(1+e^{-k})H$ is strike-dependent: at $k=-0.3$ it is $2.35\,H$, at
$k=+0.3$ only $1.74\,H$.  The put wing reads farther along the old
curve than the call wing does.

That strike dependence is exactly what the linear transport
\eqref{eq:ssrshift} at $\mathcal{R}=2$ flattens into one uniform
level.  So the frozen-field idea now has \emph{two} closed-form
candidates: the linear transport, and the relabeling
\eqref{eq:ssrell}.  Panel~(c) of \cref{fig:ssrhalf} measures their
disagreement on the real frozen node, for a
$-\MacSsrHalfMoveWingPct\%$ move: the difference is zero at the money
--- both say $2s_0H$ there --- and grows into the wings, to
$+\MacSsrHalfGapPutBp$ vol bp at the span's put edge ($k=-0.24$) and
$-\MacSsrHalfGapCallBp$ at its call edge ($k=+0.15$).  (The
asymmetric span is not a whim: it is the widest window on which the
relabeling's read-back point $\ell(k,H)$ stays inside the node's own
quoted strikes; appendix~\ref{app:ssrprotocol}.)

Which one is the frozen field's true answer?  Here the chapter owes
the reader a confession about the derivation just given.  The
half-rule was derived by averaging along a chord in
\emph{log}-strike; the relabeling by matching midpoints in
\emph{dollar} strike.  At the money, at first order, the choice of
chart cannot matter --- both gave $2s_0H$.  In the wings the two
charts stop agreeing, and neither rule of thumb outranks the other.
The honest arbiter is the computation itself: reprice the frozen field
and see.  That is the next section --- and the answer turns out to
depend on something today's smile does not show.
